% Translated from sections/01-introduction.tex on 2026-09-23. The English file
% is canonical. Change it first, then carry the change here.
\section{Introduction}
\label{sec:introduction}

Depuis un an, je construis une chaîne de traitement pour deux œuvres liées de
l'imprimé français du \sxvi. \emph{Amadis de Gaule} est un roman de
chevalerie écrit d'abord en espagnol par Garci Rodríguez de Montalvo. Sa
version française commence en 1540 avec la traduction de Nicolas Herberay des
Essarts et connaît tout de suite un grand succès. Des traducteurs, à Paris puis
à Lyon, y ajoutent sans cesse de nouveaux livres, et en 1615 le cycle français
en compte \ocrBooks{}~\cite{debuzon2011chappuys}. Le \emph{Trésor des Amadis},
imprimé pour la première fois en 1559, rassemble les harangues, lettres,
complaintes et défis du roman jugés les plus éloquents. Il devait servir de
modèle aux lecteurs désireux d'apprendre à écrire des lettres et à bien parler,
et il a été réimprimé et augmenté jusqu'en 1606~\cite{eman2026tresor}. La
chaîne vise une transcription fidèle du cycle, et un moyen de retrouver d'où
vient chaque extrait du \emph{Trésor}.

Cet imprimé résiste à la reconnaissance optique courante. Le \qtr{ſ} long est
un caractère distinct, que les modèles modernes lisent comme un \qtr{f}. Les
premiers volumes sont composés en \emph{bâtarde}, un caractère pour lequel il
existe peu de données d'entraînement modernes. Les chapitres s'ouvrent sur des
lettrines de plusieurs lignes de haut, que les segmenteurs prennent pour des
images ou du bruit. Titres courants, réclames et signatures se trouvent dans le
bloc de texte. La foliotation est en chiffres romains dans une moitié du cycle
et en chiffres arabes dans l'autre, et les numérisations viennent de plusieurs
bibliothèques, de qualité inégale, dont l'une appose son tampon sur chaque
page. La figure~\ref{fig:plate} montre une page typique.

La chaîne comporte trois étapes:
\begin{enumerate}
  \item \textbf{Reconnaissance.} Un modèle kraken, affiné à partir de
    CATMuS-Print sur des pages du \emph{Trésor}, lit chaque page.
  \item \textbf{Correction.} Les mots dont le modèle n'était pas sûr sont
    marqués, et un modèle de langue local est chargé de les corriger. Une
    correction n'est gardée que si le reste de la ligne revient inchangé.
  \item \textbf{Localisation des passages.} Chaque extrait du \emph{Trésor} est
    recherché dans le cycle transcrit, et le localisateur renvoie le passage
    correspondant et le Livre auquel il appartient.
\end{enumerate}

Ce rapport mesure chaque étape à son tour. Pour la reconnaissance, je donne la
courbe de validation de l'affinage et compare deux entraînements
(section~\ref{sec:training}); aucun jeu de test distinct n'a été évalué. Pour
la correction, je donne ce que la règle d'acceptation garde et ce qu'elle
refuse sur un échantillon de \corrPages{}~pages (section~\ref{sec:e2}). Pour la
localisation, je donne la fréquence à laquelle le localisateur trouve le bon
Livre, vérifiée d'après le catalogue de l'éditeur (section~\ref{sec:e5}). La
section~\ref{sec:e6} ajoute le débit et la détection des chapitres. Les
expériences de correction et de localisation ont été conçues par écrit avant
d'être menées.
